\documentclass[11pt]{article}
\usepackage[utf8]{inputenc}
\usepackage[T1]{fontenc}
\usepackage[a4paper,margin=1in]{geometry}
\usepackage{lmodern}
\usepackage{microtype}
\usepackage{amsmath,amssymb,mathtools}
\usepackage{mathrsfs}
\usepackage{tikz-cd}
\usepackage{enumitem}
\usepackage{longtable,booktabs,array}
\usepackage[hidelinks]{hyperref}
\setlength{\parindent}{0pt}
\setlength{\parskip}{0.55em}
\emergencystretch=3em
\hbadness=10000
\newcommand{\K}{K}
\newcommand{\Gr}{\operatorname{Gr}}
\newcommand{\Fil}{\operatorname{Fil}}
\newcommand{\ch}{\operatorname{ch}}
\newcommand{\Todd}{\operatorname{Todd}}
\newcommand{\cl}{\operatorname{cl}}
\newcommand{\Parf}{\operatorname{Parf}}
\newcommand{\Tor}{\operatorname{Tor}}
\newcommand{\Spec}{\operatorname{Spec}}
\newcommand{\RRR}{\mathrm{RRR}}
\providecommand{\codim}{\operatorname{codim}}
\providecommand{\Inv}{\operatorname{Inv}}
\providecommand{\Pic}{\operatorname{Pic}}
\providecommand{\Z}{\mathbb Z}
\providecommand{\Q}{\mathbb Q}
\providecommand{\Ahat}{\widehat A}
\providecommand{\Ahp}{\widehat A^{+}}
\providecommand{\That}{\widehat T}
\providecommand{\rad}{\operatorname{rad}}
\providecommand{\Gl}{\operatorname{Gl}}
\providecommand{\GL}{\operatorname{GL}}
\providecommand{\ch}{\operatorname{ch}}
\providecommand{\Todd}{\operatorname{Todd}}
\providecommand{\Gr}{\operatorname{Gr}}

% Core macro additions
\providecommand{\Atil}{\widetilde A}
\providecommand{\Ctil}{\widetilde C}
\providecommand{\K}{\mathrm K}
\providecommand{\cC}{\mathcal C}
\providecommand{\cF}{\mathscr F}
\providecommand{\eps}{\varepsilon}
\providecommand{\ellc}{\ell}
\providecommand{\Tor}{\operatorname{Tor}}
\providecommand{\Coker}{\operatorname{Coker}}
\providecommand{\Ker}{\operatorname{Ker}}
\providecommand{\Supp}{\operatorname{Supp}}
\providecommand{\codim}{\operatorname{codim}}
\providecommand{\Gr}{\operatorname{G}}
\providecommand{\rank}{\operatorname{rank}}
\providecommand{\rang}{\operatorname{rang}}
\providecommand{\cO}{\mathcal O}

\hypersetup{hypertexnames=false}
\providecommand{\R}{\mathbb R}
\providecommand{\N}{\mathbb N}
\begin{document}

% Local macros for Exposé V, §§6--7
\providecommand{\Ghat}{\widehat G}
\providecommand{\Aqhat}{\widehat A_{\mathbb Q}}
\providecommand{\Ch}{\operatorname{Ch}}
\providecommand{\Log}{\operatorname{Log}}
\providecommand{\id}{\operatorname{id}}
\providecommand{\gr}{\operatorname{gr}}
\providecommand{\Hom}{\operatorname{Hom}}
\providecommand{\epsi}{\varepsilon}
\providecommand{\clgr}{\operatorname{cl}}
\providecommand{\Fil}{\operatorname{Fil}}
\providecommand{\Gr}{\operatorname{Gr}}
\providecommand{\Sym}{\operatorname{Sym}}

\begin{center}
{\Large SGA 6 — Exposé V, Anneau de Chern et opérations d'Adams}\\[0.4em]
{\large Pages 354--371 du scan original}
\end{center}


\section*{Exposé V, 6. Anneau de Chern : suite}

\textbf{Définition 6.2.} Le foncteur défini en 6.1 sur la catégorie des \(K\)-algèbres graduées, à valeurs dans la catégorie des \(K\)-\(\lambda\)-algèbres, est appelé \emph{anneau de Chern}; sa valeur pour \(A\in\operatorname{Ob}(\mathscr C)\) est notée \(\Ch(A)\), ou, s'il y a des risques de confusion, \(\Ch_K(A)\).

Des formules (6.1.1), (6.1.2), (6.1.4), et (3.6.2), on déduit les relations
\begin{equation}\tag{6.2.1}
 (1,1+X)(1,1+Y)=(1,1+X+Y),
 \qquad
 \lambda^i(1,1+X)=0\quad\text{pour }i\ge2,
\end{equation}
relations qui servaient à définir l'anneau de Chern dans [1]. À partir de ces formules, on peut inversement retrouver les relations (6.1.1) et (6.1.4).

Indiquons la raison pour laquelle on introduit l'anneau de Chern. On se place dans la situation de l'exemple 3.9, et on désigne par \(A\) le gradué associé, pour la filtration définie en 3.10, au \(K^\bullet\) d'un topos commutativement localement annelé, que l'on suppose être un \(\lambda\)-anneau pour simplifier. \(K\) est l'anneau d'augmentation. On verra plus loin que l'on peut définir pour tout élément d'un \(\lambda\)-anneau augmenté des classes de Chern à valeurs dans le gradué associé (6.7). À tout faisceau localement libre de type fini, on peut donc associer un élément de \(\Ch(A)\), dont la composante sur le facteur \(K\) est le rang du faisceau, et la composante sur le facteur \(1+\widehat A^+\) la somme des classes de Chern. Les formules (6.2.1) sont alors les formules usuelles qui donnent les classes de Chern et le rang du produit tensoriel de deux faisceaux inversibles et des puissances extérieures d'un faisceau inversible, c'est-à-dire que l'application que l'on vient de définir de \(K^\bullet\) dans \(\Ch(A)\) est un \(\lambda\)-homomorphisme.

\textbf{6.3.} Pour toute \(K\)-algèbre graduée \(A\), on désigne par \(A_{\mathbb Q}\) la \(K\)-algèbre \(A\otimes\mathbb Q\), graduée de façon évidente. D'après 1.13, on définit un homomorphisme fonctoriel
\[
  1+\widehat A^+\xrightarrow{\ \ch\ }\widehat A_{\mathbb Q}^{+}
\]
en donnant sa valeur pour \(1+X\), soit
\begin{equation}\tag{6.3.1}
  \ch(1+X)=\exp(X)-1.
\end{equation}
Cet homomorphisme est compatible avec la multiplication définie en 6.1 sur \(1+\widehat A^+\). Il suffit, d'après 1.13, de le montrer pour les éléments \(1+X,1+Y\) de l'anneau universel \(1+K[X,Y]^+\). On a
\[
\begin{aligned}
 \ch\bigl((1+X)*(1+Y)\bigr)
 &= \ch\!\left(\frac{1+X+Y}{(1+X)(1+Y)}\right)\\
 &= \ch(1+X+Y)-\ch(1+X)-\ch(1+Y)\\
 &= \exp(X+Y)-\exp(X)-\exp(Y)+1\\
 &= \ch(1+X)\,\ch(1+Y).
\end{aligned}
\]

Notons \(\eta\) l'endomorphisme, fonctoriel par rapport à \(A\), de \(\widehat A_{\mathbb Q}^{+}\), défini par
\[
 \eta\!\left(\sum_{k\ge1} a_k t^k\right)
 =
 \sum_{k\ge1}\frac{(-1)^{k-1}}{(k-1)!}\,a_k t^k.
\]
Alors, pour tout \(\varphi\in1+\widehat A^+\),
\begin{equation}\tag{6.3.2}
  \ch(\varphi)=\eta(\Log\varphi).
\end{equation}
En effet, \(\ch\) et \(\eta\circ\Log\) sont deux homomorphismes fonctoriels de groupes abéliens, et il suffit qu'ils soient égaux lorsqu'ils coïncident sur \(1+X\), ce qui est évident.

On en déduit que \(\ch\) commute aux opérations de \(K\). En effet, il suffit encore de le montrer pour \(1+X\), et, si \(a\in K\),
\[
\begin{aligned}
 \ch((1+X)^a)
 &=\ch(\exp(a\Log(1+X)))\\
 &=\eta(a\Log(1+X))\\
 &=a\,\eta(\Log(1+X))\\
 &=a\,\ch(1+X).
\end{aligned}
\]
On peut donc de façon évidente étendre \(\ch\) en un homomorphisme de \(K\)-algèbres
\[
  K\times(1+\widehat A^+)\longrightarrow K\oplus\widehat A_{\mathbb Q}^{+};
\]
on a donc
\begin{equation}\tag{6.3.3}
  \ch(a,\varphi)=a+\eta(\Log\varphi).
\end{equation}

\textbf{Définition 6.4.} L'homomorphisme
\[
  \ch:\Ch A\longrightarrow K\oplus\widehat A_{\mathbb Q}^{+}
\]
qui vient d'être défini s'appelle le \emph{caractère de Chern}.

Si \(A\) est de caractéristique \(0\), c'est-à-dire est une \(\mathbb Q\)-algèbre, le caractère de Chern est un isomorphisme. En effet, soit \(\psi\in K\oplus\widehat A^+\); on peut écrire \(\psi=a+\psi_1\), où \(\psi_1\) est sans terme de degré nul. On définit alors un morphisme inverse par
\[
  \psi\longmapsto (a,\exp(\eta^{-1}(\psi_1))).
\]

\textbf{Proposition 6.5.} Soit \(A\) une \(K\)-algèbre graduée. Considérons sur \(\Ch(A)\) la filtration définie comme suit, la filtration par le degré :
\begin{equation}\tag{6.5.1}
  (\Ch(A))_0=\Ch(A),\qquad
  (\Ch(A))_n=0\times(1+\widehat A^{\ge n})\quad(n\ge1),
\end{equation}
où \(\widehat A^{\ge n}=\prod_{i\ge n}A^i\). On a les congruences suivantes.

\begin{enumerate}[label=\roman*)]
\item Soient \(n\ge1\) et \(\varphi\in(\Ch A)_n\). Posons \(\varphi=1+a_n+\cdots\). Alors
\begin{equation}\tag{6.5.2}
  \ch(\varphi)=\frac{(-1)^{n-1}}{(n-1)!}\,a_n+\cdots .
\end{equation}

\item Soient \(\varphi\in(\Ch A)_m\), \(\psi\in(\Ch A)_n\), et posons
\(\varphi=1+a_m+\cdots\), \(\psi=1+b_n+\cdots\). Alors
\begin{equation}\tag{6.5.3}
  \varphi*\psi
  =
  1-\frac{(m+n-1)!}{(m-1)!(n-1)!}\,a_m b_n+\cdots .
\end{equation}
Par suite, la filtration par le degré est une filtration d'anneau.
\end{enumerate}

i) Le caractère de Chern est par construction un morphisme de foncteurs filtrés. De la congruence \(\varphi\equiv1+a_n\pmod{(\Ch A)_{n+1}}\), on déduit
\[
  \ch(\varphi)=\ch(1+a_n)\pmod{\widehat A^{\ge n+1}},
\]
or \(\ch(1+a_n)=\eta(\Log(1+a_n))\), donc
\[
  \ch(\varphi)=\frac{(-1)^{n-1}}{(n-1)!}a_n+\cdots .
\]

ii) Par substitution, il suffit de montrer la formule (6.5.3) pour
\[
 A=K[X_i,Y_j]_{i,j\in\mathbb N}
\]
et pour \(\varphi=1+\sum_{k\ge m}X_k\), \(\psi=1+\sum_{k\ge n}Y_k\). Comme \(K\) est sans torsion dans le cas universel, on peut plonger cet anneau dans son tensorisé par \(\mathbb Q\), et on voit ainsi que \(\ch\) est injectif. Posons \(\varphi*\psi=1+Z_p+\cdots\), \(Z_p\) étant le premier terme non nul et différent de \(1\) de \(\varphi*\psi\). La multiplicativité de \(\ch\) nous donne, d'après (6.5.2),
\[
  1+\frac{(-1)^{p-1}}{(p-1)!}Z_p+\cdots
  =
  \left(1+\frac{(-1)^{m-1}}{(m-1)!}X_m+\cdots\right)
  \left(1+\frac{(-1)^{n-1}}{(n-1)!}Y_n+\cdots\right),
\]
d'où la valeur de \(Z_p\) en identifiant les termes de plus bas degré.

\textbf{6.6.} Le \(\lambda\)-anneau \(\Ch A\) est augmenté vers \(K\). On définit en effet un \(\lambda\)-homomorphisme de \(\Ch A\) dans \(K\) par \((a,\varphi)\mapsto a\). Il est donc muni d'une filtration canonique de \(\lambda\)-anneau augmenté. Nous allons la comparer à celle qui vient d'être définie.

\textbf{Proposition 6.6.1.} Soit \(x\in\Ch A\) un élément de la forme \(x=(n,1+\cdots+a_n)\). Alors :
\[
\begin{array}{ll}
\text{i)}& \lambda^i(x)=0\quad\text{pour }i>n,\\[0.3em]
\text{ii)}& \lambda^n(x)=(1,1+a_1),\\[0.3em]
\text{iii)}& \displaystyle\sum_{i=0}^n(-1)^i\lambda^i(x)
       =(0,1-(n-1)!a_n+\cdots),
\end{array}
\]
les termes non écrits étant de degré \(>n\).

Par fonctorialité, il suffit de faire la démonstration dans le cas où \(A=K[X_1,\ldots,X_n]\), filtré par le poids par rapport aux \(X_i\), \(X_i\) étant de poids \(i\), et pour l'élément \((n,1+X_1+\cdots+X_n)\). On plonge \(A\) dans \(A'=K[T_1,\ldots,T_n]\) par les fonctions symétriques, et on obtient ainsi un homomorphisme injectif. Il vient alors
\[
 (n,1+X_1+\cdots+X_n)
 =
 \left(n,\prod_{i=1}^n(1+T_i)\right)
 =
 \sum_{i=1}^n(1,1+T_i).
\]
D'après (6.2.1), on en déduit
\[
  \lambda_t(x)=\prod_{i=1}^n\bigl(1+(1,1+T_i)t\bigr),
\]
et on en tire
\[
  \lambda^k(x)=0\quad(k>n),\qquad
  \lambda^n(x)=\prod_{i=1}^n(1,1+T_i)
   =\left(1,1+\sum_{i=1}^nT_i\right)=(1,1+X_1).
\]
Enfin
\[
  \sum_{i=1}^n(-1)^i\lambda^i(x)=\lambda_{-1}(x)
  =
  \prod_{i=1}^n\bigl(1-(1,1+T_i)\bigr)
  =
  (-1)^n\prod_{i=1}^n(0,1+T_i).
\]
Or, d'après (6.5.3),
\[
  \prod_{i=1}^n(0,1+T_i)
  =
  \left(0,1+(-1)^{n-1}(n-1)!\prod_{i=1}^nT_i+\cdots\right),
\]
soit
\[
  \lambda_{-1}(x)=(0,1-(n-1)!X_n+\cdots).
\]

\textbf{Corollaire 6.6.2.} La relation 6.6.1(iii) est vraie pour tout élément \(X\) d'augmentation \(n\).

Pour le montrer, nous utiliserons le lemme suivant.

\textbf{Lemme 6.6.3.} Pour tout \(n\), \((\Ch A)_n\) est un \(\lambda\)-idéal de \(\Ch A\).

On peut supposer \(n>1\). Soit \(x=(0,\varphi)\in(\Ch A)_n\). On peut écrire
\[
  \lambda^i(0,\varphi)=(0,1+Q^i_{i,1}+\cdots+Q^i_{i,j}+\cdots),
\]
les \(Q^i_{i,j}\) étant les composantes isobares de poids \(j\) des \(Q^i_i\) de (3.6). Pour tout \(j\), les polynômes \(Q^i_{i,j}\) sont des polynômes par rapport aux \(j\) premiers coefficients de \(\varphi\). Si \(j<n\), ils sont nuls; par suite les \(\lambda^i(x)\) appartiennent à \((\Ch A)_n\), et \((\Ch A)_n\) est un \(\lambda\)-idéal.

Soit alors \(x=(n,1+a_1+\cdots+a_n+\cdots)\) un élément d'augmentation \(n\) de \(\Ch A\). On a
\[
  x\equiv (n,1+\cdots+a_n)\pmod{(\Ch A)_{n+1}},
\]
et par suite
\[
  \lambda_{-1}(x)\equiv\lambda_{-1}(n,1+\cdots+a_n)
  =(0,1-(n-1)!a_n+\cdots),
\]
d'où le corollaire.

\textbf{Corollaire 6.6.4.} Soit \(\epsilon\) l'homomorphisme d'augmentation de \(\Ch A\). Alors, pour tout \(x\in\Ch A\), on a
\[
  \gamma^n(x-\epsilon(x))
  =
  (0,1+(-1)^{n-1}(n-1)!a_n+\cdots).
\]
Par définition,
\[
  \gamma^n(x-\epsilon(x))
  =
  (-1)^n\sum_{i=0}^n(-1)^i\lambda^i(x+n-\epsilon(x)).
\]
Comme \(x+n-\epsilon(x)\) est d'augmentation \(n\), on peut lui appliquer le corollaire 6.6.2, d'où le résultat.

\textbf{Corollaire 6.6.5.} Avec les notations de 3.10,
\[
  \Fil^n\Ch A\subset(\Ch A)_n.
\]
D'après le corollaire précédent, \(\gamma^n(x-\epsilon(x))\in(\Ch A)_n\). Le corollaire résulte alors de ce que les deux filtrations sont des filtrations d'anneau, et de ce que les \(\gamma^n(x-\epsilon(x))\), pour \(n\) et \(x\) variables, engendrent la filtration de \(\lambda\)-anneau.

\textbf{Corollaire 6.6.7.} Soit \(A\) une \(K\)-algèbre graduée de caractéristique nulle, à degrés bornés. Alors les deux filtrations définies sur \(\Ch A\) sont identiques.

Supposons que \(\Fil^pA=0\), et soit \(x=(0,1+a_n+\cdots+a_p)\) un élément de \((\Ch A)_n\). Posons
\[
  x'=\left(0,1+\frac{(-1)^{n-1}}{(n-1)!}a_n+a_{n+1}+\cdots+a_p\right).
\]
On a alors
\[
  \gamma^n(x')=(0,1+a_n+\cdots+a'_p),
\]
car \(x'\) est d'augmentation nulle. Par suite \(x-\gamma^n(x')\in(\Ch A)_{n+1}\). En itérant l'opération \(p-n\) fois, on arrive à \((\Ch A)_{p+1}\), qui est nul, et on a donc écrit \(x\) comme un élément de \(\Fil^n\Ch A\). Compte tenu de (6.6.5), les deux filtrations sont donc identiques.

\textbf{6.7.} On suppose maintenant que \(A\) est un \(\lambda\)-anneau augmenté vers \(K\), l'homomorphisme d'augmentation étant noté \(\epsilon\). On dispose donc sur \(A\) de la filtration canonique de \(\lambda\)-anneau augmenté, et on peut donc associer à \(\Gr A\) son anneau de Chern.

Si \(x\in A\), alors \(\gamma_t(x-\epsilon(x))\in\Ghat(A)\) par définition de la filtration de \(\lambda\)-anneau augmenté. On note \(c(x)\), et on appelle \emph{classe de Chern totale} de \(x\), l'image de \(\gamma_t(x-\epsilon(x))\) dans
\[
  \Ghat(A)/\Ghat'(A)=1+\widehat{\Gr A}^{+}.
\]
On note \(c^i(x)\), et on appelle \(i\)-ième classe de Chern de \(x\), le \(i\)-ième coefficient de \(c(x)\). On a donc
\begin{equation}\tag{6.7.1}
  c^i(x)=\cl_{\Gr^iA}\!\left(\gamma^i(x-\epsilon(x))\right),
  \qquad
  c(x)=1+c^1(x)+\cdots .
\end{equation}
On pose alors
\begin{equation}\tag{6.7.2}
  \widetilde c(x)=(\epsilon(x),c(x)),
\end{equation}
et on appelle \(\widetilde c(x)\) la classe de Chern complétée de \(x\).

\textbf{Proposition 6.8.} La classe de Chern complétée
\[
  \widetilde c:A\longrightarrow\Ch(\Gr A)
\]
est un \(K\)-\(\lambda\)-homomorphisme de \(K\)-\(\lambda\)-algèbres augmentées.

Il est clair que l'homomorphisme ainsi défini est additif. Pour vérifier qu'il est multiplicatif, on écrit
\[
\begin{aligned}
\gamma_t(xy-\epsilon(xy))
&=\gamma_t\bigl(\epsilon(y)(x-\epsilon(x))\bigr)
  \gamma_t\bigl(\epsilon(x)(y-\epsilon(y))\bigr)
  \gamma_t\bigl((x-\epsilon(x))(y-\epsilon(y))\bigr)\\
&=\bigl(\gamma_t(x-\epsilon(x))\bigr)^{\epsilon(y)}
  \bigl(\gamma_t(y-\epsilon(y))\bigr)^{\epsilon(x)}
  \bigl(\gamma_t(x-\epsilon(x))\circ\gamma_t(y-\epsilon(y))\bigr),
\end{aligned}
\]
d'où le résultat d'après (6.1.2), en passant au quotient. Si \(a\in K\), \(\widetilde c(a)=(a,1)\); donc \(\widetilde c\) est un homomorphisme de \(K\)-algèbres, compatible en outre aux augmentations.

Il reste enfin à voir que c'est un \(\lambda\)-homomorphisme. Il faut donc montrer la commutativité du diagramme
\[
\begin{tikzcd}[column sep=large,row sep=large]
A \arrow[r,"\gamma_t"] \arrow[d,"\widetilde c"'] &
\Ghat(A) \arrow[d,"\Ghat(\widetilde c)"]\\
\Ch A \arrow[r,"\gamma_t"'] &
\Ghat(\Ch A).
\end{tikzcd}
\]
Par additivité, il suffit de le montrer pour les éléments de \(K\), pour lesquels c'est évident, et pour les éléments d'augmentation nulle de \(A\), pour lesquels il résulte, d'après la définition de \(c\) et (6.1.4), de celle du diagramme (3.7.2).

\textbf{Proposition 6.9.} Soient \(A\) une \(K\)-\(\lambda\)-algèbre augmentée, et \(z\in\Fil^mA\). Alors
\[
  \gamma^m(z)=(-1)^{m-1}(m-1)!z\pmod{\Fil^{m+1}A},
\]
c'est-à-dire
\[
  c^m(z)=\cl_{\Gr^mA}\!\left((-1)^{m-1}(m-1)!z\right).
\]

Il faut voir que ces deux éléments donnent même image dans \(\Gr^mA\), c'est-à-dire que la \(m\)-ième classe de Chern de \(z\) est l'image de \((-1)^{m-1}(m-1)!z\).

On peut écrire \(z\) comme une somme
\[
  \sum a\,\gamma^{i_1}(x_1)\cdots\gamma^{i_k}(x_k),
\]
avec \(a\in A\), \(\epsilon(x_i)=0\) pour tout \(i\), et \(i_1+\cdots+i_k\ge m\). On a alors
\[
  \widetilde c(z)=
  \sum \widetilde c(a)\,
        \widetilde c(\gamma^{i_1}(x_1))\cdots
        \widetilde c(\gamma^{i_k}(x_k)).
\]
Comme \(\widetilde c\) est un \(\lambda\)-homomorphisme, on peut écrire
\[
  \widetilde c(\gamma^i(x_j))=\gamma^i(\widetilde c(x_j)).
\]
Compte tenu de la définition de \(c\), et de (6.6.4), on obtient
\[
  \widetilde c(\gamma^{i_j}(x_j))
  =
  \left(0,1+(-1)^{i_j-1}(i_j-1)!\,\gamma^{i_j}(x_j)+\cdots\right).
\]
D'après (6.5.3), et le fait que \(\widetilde c(a)=(a,1)\), on obtient finalement
\[
\begin{aligned}
\widetilde c(z)
&=\sum
\left(0,\left(1+
(-1)^{i_1+\cdots+i_k-1}
(i_1+\cdots+i_k-1)!
\,\gamma^{i_1}(x_1)\cdots\gamma^{i_k}(x_k)
+\cdots\right)^a\right)\\
&=
\left(0,1+
\sum_{\substack{(x_1,\ldots,x_k)\\ i_1+\cdots+i_k=m}}
(-1)^{m-1}(m-1)!a\,\gamma^{i_1}(x_1)\cdots\gamma^{i_k}(x_k)
+\cdots\right)\\
&=
(0,1+(-1)^{m-1}(m-1)!z+\cdots),
\end{aligned}
\]
d'où le résultat.

\textbf{Proposition 6.10.} Soient \(K\) un anneau binomial, \(R\) une \(K\)-\(\lambda\)-algèbre augmentée, \(N\) un élément de \(R\) tel que \(\lambda^k(N)=0\) pour \(k>d\), et \(x\in\Fil^iR\). Alors
\[
  \gamma^{d+i}(N,x)-(-1)^{d+i-1}(d+i-1)!x\in\Fil^{i+1}R.
\]

On peut trouver des éléments \(z_{\alpha,j}\in R\), et \(a_\alpha\in K\), tels que
\[
  x=\sum_\alpha a_\alpha\,
       \gamma^{i_1}(x_{\alpha,1})\cdots
       \gamma^{i_k}(x_{\alpha,k}),
\]
avec \(i_1+\cdots+i_k\ge i\), et \(\epsilon(x_{\alpha,j})=0\) pour tout \(\alpha\) et tout \(j\). Soit \(R_0\) la \(K\)-\(\lambda\)-algèbre engendrée par des générateurs \(N'_0\) et \(X_{\alpha,j}\), soumis aux relations
\[
  \gamma^k(N'_0)=0\quad\text{pour }k>d.
\]
On définit un homomorphisme d'augmentation \(R_0\to K\) par \(\epsilon(N'_0)=\epsilon(X_{\alpha,j})=0\). D'après (4.9.1), \(R_0\) est l'anneau des polynômes à coefficients dans \(K\) par rapport aux indéterminées \(\gamma^k(N'_0)\), \(1\le k\le d\), et \(\gamma^p(X_{\alpha,j})\), \(p\ge1\); la \(\lambda\)-filtration de \(R_0\) est identique à sa filtration par le poids, en attribuant le poids \(k\) à \(\gamma^k(N'_0)\) et à \(\gamma^k(X_{\alpha,j})\) (4.12 et 4.15).

On pose
\[
  X_0=\sum_\alpha a_\alpha\,
       \gamma^{i_1}(X_{\alpha,1})\cdots
       \gamma^{i_k}(X_{\alpha,k}).
\]
On va montrer 6.10 d'abord pour \(X_0\) et \(N'_0=N_0+d\), qui vérifient bien les hypothèses de 6.10. Compte tenu de la remarque faite plus haut sur la filtration de \(R_0\), il suffit pour cela de montrer que
\[
  \gamma^{d+i}(N_0,X_0)\gamma^d(N'_0)
  -
  (-1)^{d+i-1}(d+i-1)!X_0\gamma^d(N'_0)
  \in\Fil^{d+i+1}R_0 .
\]
Or \(\lambda_{-1}(N_0)=(-1)^d\gamma^d(N'_0)\); cette relation s'écrit donc
\[
  (-1)^d\gamma^{d+i}(X_0,(-1)^d\gamma^d(N'_0))
  -
  (-1)^{d+i-1}(d+i-1)!X_0\gamma^d(N'_0)
  \in\Fil^{d+i+1}R_0 .
\]
Posons \(z=(-1)^dX_0\gamma^d(N'_0)\). Alors \(z\in\Fil^{d+i}R_0\), et on est ramené à montrer
\[
  \gamma^{d+i}(z)-(-1)^{d+i-1}(d+i-1)!z\in\Fil^{d+i+1}R_0,
\]
ce qui résulte de 6.9.

Pour déduire 6.10 de ce cas particulier, on considère le \(K\)-\(\lambda\)-homomorphisme \(R_0\to R\) qui envoie \(N'_0\) sur \(N-d\) et \(X_{\alpha,j}\) sur \(x_{\alpha,j}\), pour tout \(\alpha\) et tout \(j\). Cet homomorphisme est compatible aux \(\lambda\)-filtrations, et envoie \(N_0\) sur \(N\) et \(X_0\) sur \(x\), d'où le résultat.

\textbf{Proposition 6.11.} On suppose que \(K\) est de caractéristique nulle. Soit \(n\ge0\) un entier. Alors le foncteur \(\Ch\) est une équivalence de catégories de la catégorie des \(K\)-algèbres graduées \(A\) telles que \(A_0=K\), \(A_k=0\) pour \(k>n\), dans la catégorie des \(K\)-\(\lambda\)-algèbres augmentées \(\Lambda\) telles que \(\Fil^{n+1}\Lambda=0\). Un foncteur quasi-inverse est donné par \(\Lambda\mapsto\Gr\Lambda\).

Soit \(A\) une \(K\)-algèbre graduée telle que \(A_k=0\) pour \(k>n\), munie de la filtration associée à sa graduation. Alors \((\Ch A)_{n+1}=0\), donc \(\Fil^{n+1}\Ch A=0\) par (6.6.5). Réciproquement, si \(\Lambda\) est une \(K\)-\(\lambda\)-algèbre augmentée telle que \(\Fil^{n+1}\Lambda=0\), alors \(\Gr^0\Lambda=K\), \(\Gr^k\Lambda=0\) pour \(k>n\).

On va définir des isomorphismes de foncteurs
\[
  A\longrightarrow \Gr(\Ch A),
  \qquad
  \Lambda\longrightarrow \Ch(\Gr\Lambda).
\]
Comme \(A\) est à graduation finie, les gradués associés aux filtrations sur \(\Ch A\) de \(\lambda\)-anneau et par le degré sont identiques (6.6.7). On définit
\[
  u:A\longrightarrow \Gr_{\deg}(\Ch A)
\]
par l'identité sur \(A_0\) et, pour \(i\ge1\), \(a_i\in A_i\), par
\[
  u(a_i)=\cl_{\Gr^i}\left(0,1+(-1)^{i-1}(i-1)!a_i\right),
\]
et on prolonge par additivité. D'après la relation (6.5.3), c'est un homomorphisme d'anneaux gradués, compatible aux opérations de \(K\). Le morphisme inverse \(u^{-1}\) peut être défini ainsi : si \(i\ge1\), \(\bar\varphi\in\Gr^i(\Ch A)\), le \(i\)-ième coefficient \(a_i\) de \(\varphi\) ne dépend pas de la façon dont \(\bar\varphi\) a été remonté dans \((\Ch A)_i\), et on pose
\[
  u^{-1}(\bar\varphi)=\frac{(-1)^{i-1}}{(i-1)!}a_i.
\]
L'homomorphisme \(u\) est donc un isomorphisme.

L'isomorphisme \(\Lambda\to\Ch(\Gr\Lambda)\) est défini par la classe de Chern complétée \(\widetilde c\). Pour vérifier qu'elle définit bien un isomorphisme, il suffit de vérifier que l'homomorphisme obtenu par passage aux gradués associés en est un, les filtrations étant discrètes. On obtient donc
\[
  \Gr\Lambda
   \xrightarrow{\Gr(\widetilde c)}
  \Gr_{\lambda}(\Ch(\Gr\Lambda))
   \simeq
  \Gr_{\deg}(\Ch(\Gr\Lambda))
   \xrightarrow{u^{-1}}
  \Gr\Lambda .
\]
Il suffit de montrer que le composé est l'identité. Comme les homomorphismes sont des homomorphismes de \(K\)-algèbres graduées, il suffit, d'après la définition de la filtration de \(\lambda\)-anneau de \(\Lambda\), de montrer que c'est l'identité pour les \(\gamma^i(x)\), avec \(\epsilon(x)=0\). On obtient alors
\[
  \Gr(\widetilde c)\bigl(\gamma^i(x)\bigr)
  =
  \cl_{\Gr^i}\bigl(\gamma^i(c(x))\bigr).
\]
Or \(\gamma^i(\widetilde c(x))=(0,1+(-1)^{i-1}(i-1)!c^i(x)+\cdots)\), d'après (6.6.4); donc
\[
  \cl_{\Gr^i}\bigl(\gamma^i(\widetilde c(x))\bigr)=u(c^i(x)),
\]
ce qui montre que \(u^{-1}\circ\Gr(\widetilde c)=\id\).

\textbf{Corollaire 6.12.} Si \(K\) est de caractéristique nulle, et si \(\Lambda\) est une \(K\)-\(\lambda\)-algèbre de filtration discrète, alors la classe de Chern complétée
\[
  \widetilde c:\Lambda\longrightarrow \Ch(\Gr\Lambda)
\]
est un isomorphisme.

La démonstration est contenue dans celle de 6.11.

\textbf{Définition 6.13.} Soit \(A\) une \(K\)-algèbre graduée, et \(\Lambda\) une \(K\)-\(\lambda\)-algèbre augmentée. On appelle \emph{théorie de Chern de \(\Lambda\) à valeurs dans \(A\)} tout \(\lambda\)-homomorphisme
\[
  \Lambda\longrightarrow \Ch(A)
\]
compatible avec les augmentations.

On définit un foncteur sur la catégorie des \(K\)-algèbres graduées, à valeurs dans la catégorie des ensembles, en associant à toute \(K\)-algèbre graduée \(A\) l'ensemble des théories de Chern de \(\Lambda\) à valeurs dans \(A\), soit
\[
  \Hom_\lambda(\Lambda,\Ch(A)).
\]

\textbf{Proposition 6.14.} Le foncteur défini en 6.13 est représentable. Si \(K\) est de caractéristique nulle, et \(\Lambda\) de filtration discrète, il est représenté par le gradué associé à \(\Lambda\).

Montrons d'abord le premier point. On considère la \(K\)-algèbre graduée engendrée par des générateurs \(c^i(x)\), avec \(i\ge1\), \(x\in\Fil^1\Lambda\), \(c^i(x)\) étant de degré \(i\), et on fait son quotient par les relations suivantes :
\[
\begin{array}{ll}
\text{i)}&
 c^k(x+y)=\displaystyle\sum_{i+j=k}c^i(x)c^j(y),
 \quad k\ge1,\ x,y\in\Fil^1\Lambda,\\[0.8em]
\text{ii)}&
 c^k(xy)=P''_k(c^i(x),c^j(y))
 \quad(\text{cf. }6.1),
 \quad k\ge1,\ x,y\in\Fil^1\Lambda,\\[0.4em]
\text{iii)}&
 c^k(ax)=R_{k,a}(c^i(x))
 \quad(\text{cf. }6.1),
 \quad k\ge1,\ a\in K,\ x\in\Fil^1\Lambda,\\[0.4em]
\text{iv)}&
 c^k(\gamma^j(x))=Q''_{j,k}(c^i(x))
 \quad(\text{cf. }6.1),
 \quad j,k\ge1,\ x\in\Fil^1\Lambda.
\end{array}
\]
Ces relations étant isobares, le quotient est une \(K\)-algèbre graduée \(\Omega\). On définit un morphisme
\[
  \widetilde c_0:\Lambda\longrightarrow \Ch(\Omega)
\]
par
\[
  x\longmapsto
  \bigl(\epsilon(x),1+c^1(x-\epsilon(x))+\cdots\bigr).
\]
Les relations imposées montrent que c'est une théorie de Chern de \(\Lambda\) à valeurs dans \(\Omega\). Si de plus on a une théorie de Chern \(f\), à valeurs dans une \(K\)-algèbre graduée \(A\), on a de façon unique un homomorphisme \(\Omega\to A\) tel que \(f\) se factorise par \(\Ch(\Omega)\); il est défini en envoyant, pour tout \(i\ge1\), et tout \(x\in\Fil^1\Lambda\), l'élément \(c^i(x)\) sur le \(i\)-ième coefficient de \(f(x)\). Ceci est bien défini car \(f\) est un \(\lambda\)-homomorphisme.

Supposons maintenant \(K\) de caractéristique nulle, et \(\Lambda\) de filtration discrète. Soit \(n\) tel que \(\Fil^n\Lambda=0\). Alors l'algèbre graduée \(\Omega\) est nulle en degrés \(>n\). Pour le montrer, il faut montrer que, pour \(x_1,\ldots,x_p\in\Fil^1\Lambda\), et \(k_1+\cdots+k_p\ge n\), on a
\[
  c^{k_1}(x_1)\cdots c^{k_p}(x_p)=0.
\]
Or, par hypothèse,
\[
  \gamma^{k_1}(x_1)\cdots\gamma^{k_p}(x_p)=0.
\]
En appliquant le \(\lambda\)-homomorphisme \(\widetilde c_0\), on obtient
\[
  \gamma^{k_1}(\widetilde c_0(x_1))\cdots
  \gamma^{k_p}(\widetilde c_0(x_p))=(0,1).
\]
Or, d'après (6.6.4) et (6.5.3), ceci entraîne
\[
  (k_1+\cdots+k_p)!\,
  c^{k_1}(x_1)\cdots c^{k_p}(x_p)=0.
\]
Comme \(K\) est de caractéristique nulle, on a le résultat voulu.

La propriété universelle de \(\Omega\), appliquée à l'homomorphisme canonique que \(\widetilde c\) donne, donne un diagramme commutatif
\[
\begin{tikzcd}[column sep=large,row sep=large]
\Lambda \arrow[r,"\widetilde c"] \arrow[d,"\widetilde c_0"'] &
\Ch(\Gr\Lambda)\\
\Ch(\Omega) \arrow[ur] &
\end{tikzcd}
\]
D'après la proposition 6.11, l'homomorphisme \(\Omega\to\Gr\Lambda\) est un isomorphisme si et seulement si l'homomorphisme \(\Ch(\Omega)\to\Ch(\Gr\Lambda)\) en est un. Comme il est visiblement surjectif, ce dernier l'est aussi, et il suffit donc de montrer que \(\Ch(\Omega)\to\Ch(\Gr\Lambda)\) est injectif. Par fonctorialité de \(\Ch\), on obtient le diagramme
\[
\begin{tikzcd}[column sep=large,row sep=large]
\Lambda \arrow[r,"\widetilde c_0"] \arrow[d,"\widetilde c"'] &
\Ch(\Omega) \arrow[d,"\widetilde c"]\\
\Ch(\Gr\Lambda) \arrow[r,"\Ch(\widetilde c_0)"'] &
\Ch\bigl(\Gr_\Lambda(\Ch(\Omega))\bigr).
\end{tikzcd}
\]
Comme \(\Omega\) est à degrés bornés, la flèche verticale de droite est un isomorphisme par (6.12), et, en utilisant 6.11, on obtient donc un morphisme
\[
  \Gr\Lambda\longrightarrow\Omega
\]
dont le composé des deux est l'identité d'après la propriété universelle de \(\Omega\).

\section*{7. Appendice : les opérations \(\psi^k\) d'Adams}

Nous suivons ici une présentation due à J.-P. Serre.

\textbf{7.1.} Soit \(K\) un pré-\(\lambda\)-anneau. Pour tout \(x\in K\), on désigne par
\[
  \frac{d}{dt}(\lambda_t(x))
\]
la dérivée par rapport à \(t\) de \(\lambda_t(x)\). On définit alors des éléments de \(K\), notés \(\psi^k(x)\), en posant
\begin{equation}\tag{7.1.1}
  \frac{d}{dt}(\lambda_t(x))/\lambda_t(x)
  =
  \sum_{k=1}^{\infty}(-1)^{k-1}\psi^k(x)t^{k-1}.
\end{equation}
Les applications \(\psi^k:K\to K\) sont additives. En effet, on a
\[
\begin{aligned}
 \frac{d}{dt}(\lambda_t(x+y))/\lambda_t(x+y)
 &=
 \frac{d}{dt}(\lambda_t(x)\lambda_t(y))/\lambda_t(x)\lambda_t(y)\\
 &=
 \frac{d}{dt}(\lambda_t(x))/\lambda_t(x)
 +\frac{d}{dt}(\lambda_t(y))/\lambda_t(y),
\end{aligned}
\]
et par suite, pour tout \(k\),
\[
  \psi^k(x+y)=\psi^k(x)+\psi^k(y).
\]
D'après (7.1.1), les \(\psi^k\) sont des polynômes isobares de poids \(k\) par rapport aux \(\lambda^i\), avec \(i=1,\ldots,k\). Donnons par exemple les valeurs de \(\psi^1\) et \(\psi^2\) :
\[
  \psi^1(x)=\lambda^1(x)=x,
  \qquad
  \psi^2(x)=x^2-2\lambda^2(x).
\]

\textbf{7.2.} Soit \(K\) un anneau de caractéristique \(0\). Supposons données des applications \(\psi^k\) de \(K\) dans \(K\), additives, et telles que \(\psi^1=\id_K\). On définit alors une pré-\(\lambda\)-structure sur \(K\) en posant
\begin{equation}\tag{7.2.1}
  \lambda_t(x)=
  \exp\left(\sum_{k=1}^{\infty}(-1)^{k-1}\psi^k(x)t^k/k\right).
\end{equation}
La multiplicativité résulte de l'additivité des \(\psi^k\). Il y a donc équivalence entre la donnée d'une pré-\(\lambda\)-structure sur \(K\) et la donnée des opérations \(\psi^k\) correspondantes.

\textbf{7.3.} Il est clair d'après la formule (7.1.1) que les \(\psi^k\) commutent aux \(\lambda\)-homomorphismes. Si \(K,K'\) sont des pré-\(\lambda\)-anneaux de caractéristique nulle, il est nécessaire et suffisant qu'un homomorphisme \(f:K\to K'\) commute aux \(\psi^k\) pour qu'il soit un \(\lambda\)-homomorphisme; cela résulte de 7.2, les \(\lambda^i\) s'exprimant comme polynômes par rapport aux \(\psi^k\).

\textbf{Exemples 7.4.} i) Pour \(K=\mathbb Z\), muni de sa structure de \(\lambda\)-anneau, on a
\[
  \frac{d}{dt}(\lambda_t(n))/\lambda_t(n)
  =
  n(1+t)^{n-1}/(1+t)^n
  =
  n\sum_{k=1}^{\infty}(-1)^{k-1}t^{k-1}.
\]
Par suite, les \(\psi^k\) sont tous égaux à l'application identique de \(\mathbb Z\).

ii) Plaçons-nous sur la catégorie \(\mathscr C_0\) des \(\mathbb Z\)-algèbres (1.8). Les opérations \(\psi^k\) définissent d'après 7.3 des endomorphismes additifs du foncteur en \(\lambda\)-anneaux \(G\). D'après 1.10, ils sont caractérisés par leur valeur pour \(1+Xt\). On a
\[
  \frac{d}{dt}(\lambda_u(1+Xt))/\lambda_u(1+Xt)
  =
  \frac{1+Xt}{1+\{1+Xt\}u}
  =
  \sum_{k=1}^{\infty}(-1)^{k-1}\{1+Xt\}^{\circ k}u^{k-1}.
\]
Par suite
\begin{equation}\tag{7.4.1}
  \psi^k(1+Xt)=1+X^k t.
\end{equation}

\textbf{Proposition 7.5.} i) Si \(K\) est un \(\lambda\)-anneau, les opérations \(\psi^k\) vérifient les identités
\begin{equation}\tag{7.5.1}
  \psi^k(1)=1,\qquad
  \psi^k(xy)=\psi^k(x)\psi^k(y),\qquad
  \psi^k\circ\psi^{k'}=\psi^{kk'}.
\end{equation}
ii) Réciproquement, si les \(\psi^k\) vérifient ces identités et si \(K\) est de caractéristique nulle, \(K\) est un \(\lambda\)-anneau.

i) Si \(K\) est un \(\lambda\)-anneau, on a \(\lambda_t(1)=1+t\), d'où
\[
  \frac{d}{dt}(\lambda_t(1))/\lambda_t(1)
  =
  \frac1{1+t}
  =
  \sum_{k=1}^{\infty}(-1)^{k-1}t^{k-1},
\]
et par suite \(\psi^k(1)=1\).

Pour montrer la seconde des relations (7.5.1), on considère le diagramme
\begin{equation}\tag{7.5.2}
\begin{tikzcd}[column sep=large,row sep=large]
K\times K \arrow[r,"\lambda_t\times\lambda_t"] \arrow[d] &
\widehat G(K)\times\widehat G(K) \arrow[r,"g^k\times g^k"] \arrow[d,"\circ"] &
K\times K \arrow[d]\\
K \arrow[r,"\lambda_t"'] &
\widehat G(K) \arrow[r,"g^k"'] &
K ,
\end{tikzcd}
\end{equation}
où les flèches verticales sont définies par la multiplication, et les applications \(g^k\) comme suit. L'application « dérivée logarithmique » est un homomorphisme de groupes de \(G(K)\) dans \(K[[t]]\); l'application qui à une série associe son \((k-1)\)-ième coefficient, multiplié par \((-1)^{k-1}\), est un homomorphisme additif de \(K[[t]]\) dans \(K\). Le morphisme \(g^k\) sera le composé des deux. Il est clair que \(g^k\circ\lambda_t=\psi^k\), d'après la définition (7.1.1).

La seconde des formules (7.5.1) exprime donc la commutativité du diagramme écrit plus haut. Si on suppose que \(K\) est un \(\lambda\)-anneau, le carré de gauche est commutatif, et il suffit de vérifier la commutativité de celui de droite. Or il s'agit d'un diagramme d'homomorphismes de foncteurs sur \(\mathscr C_0\), ayant pour source \(\widehat G\times\widehat G\). Il suffit donc, d'après 1.10, de vérifier la commutativité pour \((1+Xt,1+Yt)\). On a alors
\[
  g^k(1+Xt)=X^k,\qquad g^k(1+Yt)=Y^k,
\]
et
\[
  g^k\bigl((1+Xt)\circ(1+Yt)\bigr)
  =
  g^k(1+XYt)
  =
  (XY)^k
  =
  g^k(1+Xt)g^k(1+Yt),
\]
ce qui donne le résultat.

Pour montrer la dernière, on remarque d'abord que le diagramme
\begin{equation}\tag{7.5.3}
\begin{tikzcd}[column sep=large,row sep=large]
\widehat G_t(K) \arrow[r,"\lambda_u"] \arrow[d,"g^k"'] &
\widehat G_u\widehat G_t(K) \arrow[d,"g^k_u"]\\
K \arrow[r,"\lambda_u"'] &
\widehat G^u(K)
\end{tikzcd}
\end{equation}
est commutatif. Comme plus haut, il suffit de le voir pour \(1+Xt\), et on a
\[
  g^{k'}_u\lambda_u(1+Xt)
  =
  g^{k'}_u(1+\{1+Xt\}u)
  =
  (1+Xt)^{\circ k'}
  =
  1+X^{k'}t,
\]
puis
\[
  g^k(1+X^{k'}t)=X^{kk'}=g^k(1+Xt).
\]
Par suite
\[
  \psi^k\circ\psi^{k'}
  =
  g^k\circ g^{k'}_u\circ\lambda_u\circ\lambda_t
  =
  g^k\circ\lambda_t\circ\psi^{k'}.
\]
Par ailleurs
\[
  \psi^k\circ\psi^{k'}
  =
  g^k\circ\lambda_t\circ\psi^{k'}
  =
  \psi^k\circ\lambda_t\circ\psi^{k'}.
\]
Or si \(K\) est un \(\lambda\)-anneau, le diagramme
\begin{equation}\tag{7.5.4}
\begin{tikzcd}[column sep=large,row sep=large]
K \arrow[r,"\lambda_t"] \arrow[d,"\psi^{k'}"'] &
\widehat G(K) \arrow[d,"\psi^{k'}"]\\
K \arrow[r,"\lambda_t"'] &
\widehat G(K)
\end{tikzcd}
\end{equation}
est commutatif d'après 7.3, puisque \(\lambda_t\) est un \(\lambda\)-homomorphisme. On en déduit aussitôt les relations cherchées.

ii) Si \(K\) est de caractéristique nulle, l'application « dérivée logarithmique » est un isomorphisme de \(\widehat G(K)\) sur \(K[[t]]\), car on peut intégrer. Par suite, deux éléments de \(\widehat G(K)\) sont égaux si et seulement si, pour tout \(k\), ils ont même image par \(g^k\).

Pour montrer que \(\lambda_t\) est un homomorphisme d'anneaux, il faut montrer la commutativité du carré de gauche du diagramme (7.5.2), et, d'après la remarque précédente, il suffit de montrer que, pour tout \(k\), le diagramme (7.5.2) est commutatif, puisque le carré de droite l'est toujours. Or ceci n'est autre que les relations
\[
  \psi^k(xy)=\psi^k(x)\psi^k(y).
\]
Pour montrer que \(\lambda_t\) est un \(\lambda\)-homomorphisme, il suffit de montrer que le diagramme (7.5.4) est commutatif pour tout \(k'\), d'après 7.3. Or, d'après la remarque qui précède, cela équivaut à montrer pour tous \(k,k'\) que
\[
  g^k\circ\psi^{k'}\circ\lambda_t
  =
  g^k\circ\lambda_t\circ\psi^{k'},
\]
c'est-à-dire les relations \(\psi^{kk'}=\psi^k\circ\psi^{k'}\). Enfin, la relation \(\lambda_t(1)=1+t\) s'obtient par intégration à partir des relations \(\psi^k(1)=1\).

\subsection*{Bibliographie}
\begin{enumerate}[label={[\arabic*]}]
\item A. Grothendieck, \emph{La théorie des classes de Chern}, Bull. Soc. math. France, 86, 1958, p. 137 à 154.
\item F. Hirzebruch, \emph{Neue topologische Methoden in der algebraischen Geometrie}.
\end{enumerate}

\end{document}
